\section{Simulacro 4}
\textbf{Cuestiones.} Un choque es un proceso irreversible: la entropía aumenta y la presión total cae, mientras que la temperatura total se conserva porque el flujo es adiabático y no hay trabajo. Solo es posible un salto de supersónico a subsónico, nunca el contrario. El ángulo de Mach $\mu=\arcsin(1/M)$ corresponde a perturbaciones infinitesimales; un choque oblicuo de intensidad finita tiene $\beta>\mu$, y ambos coinciden solo en el límite de desviación nula. La corrección de Prandtl-Glauert procede de linealizar flujo subsónico con perturbaciones pequeñas; diverge en $M\to1$, donde aparecen regiones supersónicas locales y choques, y en supersónico la teoría lineal es la de Ackeret, con otra forma.

\textbf{Pitot.} A $12000\,\mathrm m$ la capa es isoterma con $T=216.65\,\mathrm K$. Se obtiene $p=19330.4\,\mathrm{Pa}$, $\rho=0.310828\,\mathrm{kg/m^3}$, $a=295.069\,\mathrm{m/s}$ y $V=590.139\,\mathrm{m/s}$. Con $M_1=2$, el choque normal da $M_2=0.57735$, $p_2/p_1=4.5$, $\rho_2/\rho_1=2.66667$, $T_2/T_1=1.6875$ y $p_{t2}/p_{t1}=0.720874$. El Pitot mide $p_{t2}=109032\,\mathrm{Pa}$. Con la fórmula subsónica se leería $M=1.78789$ en lugar de $2$: el error es grande porque esa fórmula supone compresión isentrópica completa y aquí hay una pérdida de presión total en el choque desprendido.

\textbf{Cuña.} La relación $\tan\theta=2\cot\beta\,(M_1^2\sin^2\beta-1)/[M_1^2(\gamma+\cos2\beta)+2]$ da $\theta=17.6755^\circ$. La componente normal es $M_{n1}=1.60697$, con $p_2/p_1=2.84607$, $\rho_2/\rho_1=2.04344$ y $T_2/T_1=1.39278$. Tras el choque, $M_{n2}=0.666351$ y $M_2=M_{n2}/\sin(\beta-\theta)=1.75424$, todavía supersónico: es un choque débil. Para cada $M_1$ existe una desviación máxima; por encima de ella el choque se desprende y pasa a ser curvo.

\textbf{Perfil.} A $3000\,\mathrm m$, $T=268.65\,\mathrm K$, $\rho=0.909122\,\mathrm{kg/m^3}$ y $a=328.578\,\mathrm{m/s}$, de modo que $V=197.147\,\mathrm{m/s}$. Con $\mu=1.69364\times10^{-5}\,\mathrm{Pa\,s}$ resulta $Re=1.90486\times10^{7}$. El coeficiente incompresible es $C_l=2\pi(\alpha-\alpha_{L0})=0.603142$; con $\beta_{PG}=\sqrt{1-M^2}=0.8$ se obtiene $C_l=0.753928$. La corrección es razonable a $M=0.6$ si el perfil es delgado y la incidencia pequeña, y deja de serlo cerca del Mach crítico.

\subsection*{Procedimiento y operaciones del simulacro 4}
La atmósfera ISA da el estado antes del choque. Para el choque normal se usan las relaciones de Rankine-Hugoniot en función de $M_1$; la presión de remanso medida por un Pitot supersónico sigue la fórmula de Rayleigh. En el choque oblicuo, la relación $\theta$-$\beta$-$M$ da la desviación y las relaciones normales se aplican a $M_1\sin\beta$. En el perfil se usa $\rho Vc/\mu$ y la corrección lineal subsónica.

\begin{enumerate}

\item Temperatura de la base de la capa isoterma.

\[\fitmath{T_{11}=288.15-0.0065\cdot 11000\approx 216.65\,\mathrm{K}}\]

\item Presión a 11 km por la integral de la capa troposférica.

\[\fitmath{p_{11}=101325\cdot {\left(\frac{T_{11}}{288.15}\right)}^{\frac{9.80665}{287.05287\cdot 0.0065}}}\]

\[\fitmath{p_{11}=101325\cdot {\left(\frac{(216.65)}{288.15}\right)}^{\frac{9.80665}{287.05287\cdot 0.0065}}\approx 22632.04\,\mathrm{Pa}}\]

\item Presión a 12 km con la integral isoterma.

\[\fitmath{p_{1}=p_{11}\cdot \exp\left(\frac{-9.80665\cdot \left(12000-11000\right)}{287.05287\cdot T_{11}}\right)}\]

\[\fitmath{p_{1}=(22632.0401)\cdot \exp\left(\frac{-9.80665\cdot \left(12000-11000\right)}{287.05287\cdot (216.65)}\right)\approx 19330.38\,\mathrm{Pa}}\]

\item Densidad por la ecuación de estado con la temperatura de la capa.

\[\fitmath{\rho_{1}=\frac{p_{1}}{287.05287\cdot T_{11}}}\]

\[\fitmath{\rho_{1}=\frac{(19330.38251)}{287.05287\cdot (216.65)}\approx 0.3108278\,\mathrm{kg/m^{3}}}\]

\item Velocidad del sonido local.

\[\fitmath{a_{1}=\sqrt{1.4\cdot 287.05287\cdot T_{11}}}\]

\[\fitmath{a_{1}=\sqrt{1.4\cdot 287.05287\cdot (216.65)}\approx 295.0695\,\mathrm{m/s}}\]

\item Mach de vuelo dado.

\[\fitmath{M_{1}=2\approx 2\,}\]

\item Velocidad verdadera.

\[\fitmath{V_{1}=M_{1}\cdot a_{1}}\]

\[\fitmath{V_{1}=(2)\cdot (295.0694935)\approx 590.139\,\mathrm{m/s}}\]

\item Mach detrás del choque normal.

\[\fitmath{M_{2}=\sqrt{\frac{1+0.2\cdot {M_{1}}^{2}}{1.4\cdot {M_{1}}^{2}-0.2}}}\]

\[\fitmath{M_{2}=\sqrt{\frac{1+0.2\cdot {(2)}^{2}}{1.4\cdot {(2)}^{2}-0.2}}\approx 0.5773503\,}\]

\item Salto de presión estática a través del choque.

\[\fitmath{\mathrm{pr}=1+\frac{2\cdot 1.4}{2.4}\cdot \left({M_{1}}^{2}-1\right)}\]

\[\fitmath{\mathrm{pr}=1+\frac{2\cdot 1.4}{2.4}\cdot \left({(2)}^{2}-1\right)\approx 4.5\,}\]

\item Salto de densidad por Rankine-Hugoniot.

\[\fitmath{\mathrm{rr}=\frac{2.4\cdot {M_{1}}^{2}}{0.4\cdot {M_{1}}^{2}+2}}\]

\[\fitmath{\mathrm{rr}=\frac{2.4\cdot {(2)}^{2}}{0.4\cdot {(2)}^{2}+2}\approx 2.666667\,}\]

\item Cociente de temperaturas por gas perfecto.

\[\fitmath{\mathrm{Tr}=\frac{\mathrm{pr}}{\mathrm{rr}}}\]

\[\fitmath{\mathrm{Tr}=\frac{(4.5)}{(2.666666667)}\approx 1.6875\,}\]

\item Pérdida de presión total por choque normal.

\[\fitmath{\mathrm{ptr}=\frac{{\left(\frac{2.4\cdot {M_{1}}^{2}}{0.4\cdot {M_{1}}^{2}+2}\right)}^{\frac{1.4}{0.4}}}{{\left(\frac{2.8}{2.4}\cdot {M_{1}}^{2}-\frac{0.4}{2.4}\right)}^{\frac{1}{0.4}}}}\]

\[\fitmath{\mathrm{ptr}=\frac{{\left(\frac{2.4\cdot {(2)}^{2}}{0.4\cdot {(2)}^{2}+2}\right)}^{\frac{1.4}{0.4}}}{{\left(\frac{2.8}{2.4}\cdot {(2)}^{2}-\frac{0.4}{2.4}\right)}^{\frac{1}{0.4}}}\approx 0.7208739\,}\]

\item Presión de remanso medida tras el choque (Rayleigh-Pitot).

\[\fitmath{\mathrm{pt}_{2}=\frac{p_{1}\cdot {\left(\frac{2.4}{2}\cdot {M_{1}}^{2}\right)}^{\frac{1.4}{0.4}}}{{\left(\frac{2\cdot 1.4}{2.4}\cdot {M_{1}}^{2}-\frac{0.4}{2.4}\right)}^{\frac{1}{0.4}}}}\]

\[\fitmath{\mathrm{pt}_{2}=\frac{(19330.38251)\cdot {\left(\frac{2.4}{2}\cdot {(2)}^{2}\right)}^{\frac{1.4}{0.4}}}{{\left(\frac{2\cdot 1.4}{2.4}\cdot {(2)}^{2}-\frac{0.4}{2.4}\right)}^{\frac{1}{0.4}}}\approx 109031.9\,\mathrm{Pa}}\]

\item Mach que se leería aplicando la fórmula subsónica a esa presión.

\[\fitmath{\mathrm{Mfalse}=\sqrt{5\cdot \left({\left(\frac{\mathrm{pt}_{2}}{p_{1}}\right)}^{\frac{0.4}{1.4}}-1\right)}}\]

\[\fitmath{\mathrm{Mfalse}=\sqrt{5\cdot \left({\left(\frac{(109031.8784)}{(19330.38251)}\right)}^{\frac{0.4}{1.4}}-1\right)}\approx 1.787889\,}\]

\item Ángulo del choque oblicuo en radianes.

\[\fitmath{\beta=\frac{\pi\cdot 40}{180}\approx 0.6981317\,\mathrm{rad}}\]

\item Mach aguas arriba del choque oblicuo.

\[\fitmath{\mathrm{Mi}=2.5\approx 2.5\,}\]

\item Componente normal del Mach incidente.

\[\fitmath{\mathrm{Mn}_{1}=\mathrm{Mi}\cdot \sin\left(\beta\right)}\]

\[\fitmath{\mathrm{Mn}_{1}=(2.5)\cdot \sin\left((0.6981317008)\right)\approx 1.606969\,}\]

\item Ángulo de desviación de la cuña por la relación theta-beta-M.

\[\fitmath{\theta=\arctan\left(\frac{\frac{2}{\tan\left(\beta\right)}\cdot \left({\mathrm{Mi}}^{2}\cdot {\sin\left(\beta\right)}^{2}-1\right)}{{\mathrm{Mi}}^{2}\cdot \left(1.4+\cos\left(2\cdot \beta\right)\right)+2}\right)}\]

\[\fitmath{\theta=\arctan\left(\frac{\frac{2}{\tan\left((0.6981317008)\right)}\cdot \left({(2.5)}^{2}\cdot {\sin\left((0.6981317008)\right)}^{2}-1\right)}{{(2.5)}^{2}\cdot \left(1.4+\cos\left(2\cdot (0.6981317008)\right)\right)+2}\right)\approx 0.3084949\,\mathrm{rad}}\]

\item Ángulo de la cuña en grados.

\[\fitmath{\mathrm{thetadeg}=\frac{\theta\cdot 180}{\pi}}\]

\[\fitmath{\mathrm{thetadeg}=\frac{(0.3084949364)\cdot 180}{\pi}\approx 17.67546\,{}^{\circ}}\]

\item Salto de presión por la componente normal.

\[\fitmath{\mathrm{pr}_{2}=1+\frac{2\cdot 1.4}{2.4}\cdot \left({\mathrm{Mn}_{1}}^{2}-1\right)}\]

\[\fitmath{\mathrm{pr}_{2}=1+\frac{2\cdot 1.4}{2.4}\cdot \left({(1.606969024)}^{2}-1\right)\approx 2.846074\,}\]

\item Salto de densidad por la componente normal.

\[\fitmath{\mathrm{rr}_{2}=\frac{2.4\cdot {\mathrm{Mn}_{1}}^{2}}{0.4\cdot {\mathrm{Mn}_{1}}^{2}+2}}\]

\[\fitmath{\mathrm{rr}_{2}=\frac{2.4\cdot {(1.606969024)}^{2}}{0.4\cdot {(1.606969024)}^{2}+2}\approx 2.043443\,}\]

\item Mach normal posterior.

\[\fitmath{\mathrm{Mn}_{2}=\sqrt{\frac{1+0.2\cdot {\mathrm{Mn}_{1}}^{2}}{1.4\cdot {\mathrm{Mn}_{1}}^{2}-0.2}}}\]

\[\fitmath{\mathrm{Mn}_{2}=\sqrt{\frac{1+0.2\cdot {(1.606969024)}^{2}}{1.4\cdot {(1.606969024)}^{2}-0.2}}\approx 0.6663512\,}\]

\item Mach aguas abajo, recomponiendo con la desviación.

\[\fitmath{\mathrm{Mo}_{2}=\frac{\mathrm{Mn}_{2}}{\sin\left(\beta-\theta\right)}}\]

\[\fitmath{\mathrm{Mo}_{2}=\frac{(0.6663511504)}{\sin\left((0.6981317008)-(0.3084949364)\right)}\approx 1.754237\,}\]

\item Temperatura ISA a 3000 m.

\[\fitmath{\mathrm{Tc}=288.15-0.0065\cdot 3000\approx 268.65\,\mathrm{K}}\]

\item Presión ISA a 3000 m.

\[\fitmath{\mathrm{pc}=101325\cdot {\left(\frac{\mathrm{Tc}}{288.15}\right)}^{\frac{9.80665}{287.05287\cdot 0.0065}}}\]

\[\fitmath{\mathrm{pc}=101325\cdot {\left(\frac{(268.65)}{288.15}\right)}^{\frac{9.80665}{287.05287\cdot 0.0065}}\approx 70108.53\,\mathrm{Pa}}\]

\item Densidad por gas perfecto.

\[\fitmath{\mathrm{rhoc}=\frac{\mathrm{pc}}{287.05287\cdot \mathrm{Tc}}}\]

\[\fitmath{\mathrm{rhoc}=\frac{(70108.5265)}{287.05287\cdot (268.65)}\approx 0.9091219\,\mathrm{kg/m^{3}}}\]

\item Velocidad del sonido ISA.

\[\fitmath{\mathrm{ac}=\sqrt{1.4\cdot 287.05287\cdot \mathrm{Tc}}}\]

\[\fitmath{\mathrm{ac}=\sqrt{1.4\cdot 287.05287\cdot (268.65)}\approx 328.5779\,\mathrm{m/s}}\]

\item Viscosidad de Sutherland.

\[\fitmath{\mathrm{muc}=\frac{1.716\times10^{-5}\cdot {\left(\frac{\mathrm{Tc}}{273.15}\right)}^{1.5}\cdot \left(273.15+110.4\right)}{\mathrm{Tc}+110.4}}\]

\[\fitmath{\mathrm{muc}=\frac{1.716\times10^{-5}\cdot {\left(\frac{(268.65)}{273.15}\right)}^{1.5}\cdot \left(273.15+110.4\right)}{(268.65)+110.4}\approx 1.69364\times10^{-5}\,\mathrm{Pa\,s}}\]

\item Mach de vuelo.

\[\fitmath{\mathrm{Mc}=0.6\approx 0.6\,}\]

\item Velocidad verdadera.

\[\fitmath{\mathrm{Vc}=\mathrm{Mc}\cdot \mathrm{ac}}\]

\[\fitmath{\mathrm{Vc}=(0.6)\cdot (328.5779283)\approx 197.1468\,\mathrm{m/s}}\]

\item Reynolds con cuerda de 1.8 m.

\[\fitmath{\mathrm{Re}=\frac{\mathrm{rhoc}\cdot \mathrm{Vc}\cdot 1.8}{\mathrm{muc}}}\]

\[\fitmath{\mathrm{Re}=\frac{(0.9091218612)\cdot (197.146757)\cdot 1.8}{(1.693640497\times10^{-5})}\approx 1.90486\times10^{7}\,}\]

\item Pendiente 2 pi, incidencia 4 grados, alfa0 -1.5 grados.

\[\fitmath{\mathrm{Cl}_{0}=\frac{2\cdot \pi\cdot \left(4+1.5\right)\cdot \pi}{180}\approx 0.6031425\,}\]

\item Factor de Prandtl-Glauert.

\[\fitmath{\mathrm{bPG}=\sqrt{1-{\mathrm{Mc}}^{2}}}\]

\[\fitmath{\mathrm{bPG}=\sqrt{1-{(0.6)}^{2}}\approx 0.8\,}\]

\item Corrección subsónica de Prandtl-Glauert.

\[\fitmath{\mathrm{Clc}=\frac{\mathrm{Cl}_{0}}{\mathrm{bPG}}}\]

\[\fitmath{\mathrm{Clc}=\frac{(0.6031424912)}{(0.8)}\approx 0.7539281\,}\]

\item Cociente de temperaturas por gas perfecto.

\[\fitmath{\mathrm{Tr}_{2}=\frac{\mathrm{pr}_{2}}{\mathrm{rr}_{2}}}\]

\[\fitmath{\mathrm{Tr}_{2}=\frac{(2.846074352)}{(2.043442706)}\approx 1.392784\,}\]

\end{enumerate}


\section{Simulacro 5}
\textbf{Cuestiones.} En el ciclo Otto ideal el calor cedido y el aportado son a volumen constante entre las mismas temperaturas límite, y las dos adiabáticas fijan $T_2/T_1=r^{\gamma-1}$; el rendimiento es $1-r^{1-\gamma}$ y no depende del calor aportado. En un motor real la relación está limitada por la detonación del combustible, por pérdidas de calor y por el aumento de tensiones mecánicas. En un turborreactor el rendimiento propulsivo $2V_0/(V_e+V_0)$ aumenta al acercarse $V_e$ a $V_0$, pero el empuje $\dot m(V_e-V_0)$ disminuye si no crece el gasto. El impulso específico mide empuje por unidad de peso de propulsante consumido; en vacío desaparece el término $-p_aA_e$, que resta empuje a nivel del mar.

\textbf{Otto.} Con $r=9$, $T_2=T_1r^{0.4}=722.467\,\mathrm K$ y $p_2=p_1r^{1.4}=2.05903\times10^{6}\,\mathrm{Pa}$. Con $c_v=R_a/(\gamma-1)=717.632\,\mathrm{J/(kg\,K)}$, el aporte de calor lleva a $T_3=T_2+q/c_v=2394.63\,\mathrm K$ y $p_3=p_2T_3/T_2=6.8247\times10^{6}\,\mathrm{Pa}$. La expansión termina en $T_4=T_3/r^{0.4}=994.356\,\mathrm K$. El rendimiento es $\eta=1-r^{-0.4}=0.584756$, el trabajo $w=\eta q=701708\,\mathrm{J/kg}$ y el calor cedido $q-w=498292\,\mathrm{J/kg}$, que coincide con $c_v(T_4-T_1)=498292\,\mathrm{J/kg}$.

\textbf{Turborreactor.} El empuje es $F=\dot m[(1+f)V_e-V_0]=18418.5\,\mathrm N$. La potencia propulsiva es $FV_0=4.05207\times10^{6}\,\mathrm W$ y el incremento de energía cinética $7.68974\times10^{6}\,\mathrm W$, por lo que $\eta_p=0.526945$. El consumo es $\mathrm{TSFC}=f\dot m/F=3.66479\times10^{-5}\,\mathrm{kg/(N\,s)}=0.131933\,\mathrm{kg/(N\,h)}$. Es un modelo docente: no considera presión de salida distinta de la ambiente ni sangrados.

\textbf{Cohete.} El gasto crítico es $\dot m=p_cA_t\sqrt{\gamma/(RT_c)}\,[2/(\gamma+1)]^{(\gamma+1)/[2(\gamma-1)]}=114.112\,\mathrm{kg/s}$. A la salida, $T_e=T_c(p_e/p_c)^{(\gamma-1)/\gamma}=1626.83\,\mathrm K$ y $V_e=2843.53\,\mathrm{m/s}$, con $\rho_e=0.0970566\,\mathrm{kg/m^3}$ y $A_e=\dot m/(\rho_eV_e)=0.413473\,\mathrm{m^2}$. La relación de expansión es $10.3368$ y $M_e=3.30144$. El empuje es $F_{sl}=\dot mV_e+(p_e-p_a)A_e=307393\,\mathrm N$ al nivel del mar y $F_{vac}=\dot mV_e+p_eA_e=349288\,\mathrm N$ en vacío, con $I_{sp}=274.69\,\mathrm s$ y $312.128\,\mathrm s$. Como $p_e<p_a$, la tobera está sobreexpandida al nivel del mar; con $p_e/p_a\approx0.59$ no cabe esperar desprendimiento de la capa límite, que suele aparecer con cocientes mucho menores, pero conviene comprobarlo.

\subsection*{Procedimiento y operaciones del simulacro 5}
En el ciclo Otto se encadenan las dos isentrópicas con el aporte a volumen constante. En el turborreactor se aplica la cantidad de movimiento con tobera adaptada, de modo que no hay término de presión. En la tobera de cohete se usan el gasto crítico, la velocidad de salida isentrópica y la continuidad para el área de salida; el empuje añade el término de presión $(p_e-p_a)A_e$.

\begin{enumerate}

\item Relación de compresión dada.

\[\fitmath{r=9\approx 9\,}\]

\item Temperatura de admisión.

\[\fitmath{T_{1}=300\approx 300\,\mathrm{K}}\]

\item Presión de admisión.

\[\fitmath{p_{1}=95000\approx 95000\,\mathrm{Pa}}\]

\item Calor aportado por kilogramo de mezcla.

\[\fitmath{q=1200000\approx 1200000\,\mathrm{J/kg}}\]

\item Calor específico a volumen constante de gas perfecto con gamma 1.4.

\[\fitmath{\mathrm{cv}=\frac{287.05287}{0.4}\approx 717.6322\,\mathrm{J/(kg\,K)}}\]

\item Compresión isentrópica.

\[\fitmath{T_{2}=T_{1}\cdot {r}^{0.4}}\]

\[\fitmath{T_{2}=(300)\cdot {(9)}^{0.4}\approx 722.4674\,\mathrm{K}}\]

\item Presión tras la compresión isentrópica.

\[\fitmath{p_{2}=p_{1}\cdot {r}^{1.4}}\]

\[\fitmath{p_{2}=(95000)\cdot {(9)}^{1.4}\approx 2059032\,\mathrm{Pa}}\]

\item Aporte de calor a volumen constante.

\[\fitmath{T_{3}=T_{2}+\frac{q}{\mathrm{cv}}}\]

\[\fitmath{T_{3}=(722.4674056)+\frac{(1200000)}{(717.632175)}\approx 2394.633\,\mathrm{K}}\]

\item A volumen constante la presión crece como la temperatura.

\[\fitmath{p_{3}=\frac{p_{2}\cdot T_{3}}{T_{2}}}\]

\[\fitmath{p_{3}=\frac{(2059032.106)\cdot (2394.633233)}{(722.4674056)}\approx 6824705\,\mathrm{Pa}}\]

\item Expansión isentrópica.

\[\fitmath{T_{4}=\frac{T_{3}}{{r}^{0.4}}}\]

\[\fitmath{T_{4}=\frac{(2394.633233)}{{(9)}^{0.4}}\approx 994.3562\,\mathrm{K}}\]

\item Rendimiento térmico del ciclo Otto ideal.

\[\fitmath{\eta=1-{r}^{-0.4}}\]

\[\fitmath{\eta=1-{(9)}^{-0.4}\approx 0.5847564\,}\]

\item Trabajo neto específico.

\[\fitmath{w=\eta\cdot q}\]

\[\fitmath{w=(0.5847563535)\cdot (1200000)\approx 701707.6\,\mathrm{J/kg}}\]

\item Calor cedido a volumen constante.

\[\fitmath{\mathrm{qout}=q-w}\]

\[\fitmath{\mathrm{qout}=(1200000)-(701707.6242)\approx 498292.4\,\mathrm{J/kg}}\]

\item Comprobación de calor cedido desde las temperaturas.

\[\fitmath{\mathrm{qout}_{2}=\mathrm{cv}\cdot \left(T_{4}-T_{1}\right)}\]

\[\fitmath{\mathrm{qout}_{2}=(717.632175)\cdot \left((994.356236)-(300)\right)\approx 498292.4\,\mathrm{J/kg}}\]

\item Gasto de aire.

\[\fitmath{\mathrm{mdot}=45\approx 45\,\mathrm{kg/s}}\]

\item Dosado másico combustible/aire.

\[\fitmath{f=0.015\approx 0.015\,}\]

\item Velocidad de vuelo.

\[\fitmath{V_{0}=220\approx 220\,\mathrm{m/s}}\]

\item Velocidad de salida con tobera adaptada.

\[\fitmath{\mathrm{Ve}=620\approx 620\,\mathrm{m/s}}\]

\item Empuje con tobera adaptada.

\[\fitmath{F=\mathrm{mdot}\cdot \left(\left(1+f\right)\cdot \mathrm{Ve}-V_{0}\right)}\]

\[\fitmath{F=(45)\cdot \left(\left(1+(0.015)\right)\cdot (620)-(220)\right)\approx 18418.5\,\mathrm{N}}\]

\item Potencia propulsiva.

\[\fitmath{\mathrm{Pu}=F\cdot V_{0}}\]

\[\fitmath{\mathrm{Pu}=(18418.5)\cdot (220)\approx 4052070\,\mathrm{W}}\]

\item Incremento de energía cinética de la corriente.

\[\fitmath{\mathrm{dKE}=\frac{\mathrm{mdot}\cdot \left(\left(1+f\right)\cdot {\mathrm{Ve}}^{2}-{V_{0}}^{2}\right)}{2}}\]

\[\fitmath{\mathrm{dKE}=\frac{(45)\cdot \left(\left(1+(0.015)\right)\cdot {(620)}^{2}-{(220)}^{2}\right)}{2}\approx 7689735\,\mathrm{W}}\]

\item Rendimiento propulsivo.

\[\fitmath{\eta_p=\frac{\mathrm{Pu}}{\mathrm{dKE}}}\]

\[\fitmath{\eta_p=\frac{(4052070)}{(7689735)}\approx 0.5269453\,}\]

\item Consumo específico de empuje.

\[\fitmath{\mathrm{TSFC}=\frac{f\cdot \mathrm{mdot}}{F}}\]

\[\fitmath{\mathrm{TSFC}=\frac{(0.015)\cdot (45)}{(18418.5)}\approx 3.664794\times10^{-5}\,\mathrm{kg/(N\,s)}}\]

\item Mismo consumo en unidades por hora.

\[\fitmath{\mathrm{TSFCh}=\mathrm{TSFC}\cdot 3600}\]

\[\fitmath{\mathrm{TSFCh}=(3.66479355\times10^{-5})\cdot 3600\approx 0.1319326\,\mathrm{kg/(N\,h)}}\]

\item Razón de calores específicos de los gases de combustión.

\[\fitmath{\gamma=1.2\approx 1.2\,}\]

\item Constante del gas.

\[\fitmath{R_g=380\approx 380\,\mathrm{J/(kg\,K)}}\]

\item Temperatura de cámara.

\[\fitmath{\mathrm{Tc}=3400\approx 3400\,\mathrm{K}}\]

\item Presión de cámara.

\[\fitmath{\mathrm{pc}=5000000\approx 5000000\,\mathrm{Pa}}\]

\item Área de garganta.

\[\fitmath{\mathrm{At}=0.04\approx 0.04\,\mathrm{m^{2}}}\]

\item Presión de salida de diseño.

\[\fitmath{\mathrm{pe}=60000\approx 60000\,\mathrm{Pa}}\]

\item Presión ambiente a nivel del mar.

\[\fitmath{\mathrm{pa}=101325\approx 101325\,\mathrm{Pa}}\]

\item Gasto crítico en la garganta bloqueada.

\[\fitmath{\mathrm{mdotc}=\mathrm{pc}\cdot \mathrm{At}\cdot \sqrt{\frac{\gamma}{R_g\cdot \mathrm{Tc}}}\cdot {\left(\frac{2}{\gamma+1}\right)}^{\frac{\gamma+1}{2\cdot \left(\gamma-1\right)}}}\]

\[\fitmath{\mathrm{mdotc}=(5000000)\cdot (0.04)\cdot \sqrt{\frac{(1.2)}{(380)\cdot (3400)}}\cdot {\left(\frac{2}{(1.2)+1}\right)}^{\frac{(1.2)+1}{2\cdot \left((1.2)-1\right)}}\approx 114.1115\,\mathrm{kg/s}}\]

\item Temperatura de salida isentrópica.

\[\fitmath{\mathrm{Te}=\mathrm{Tc}\cdot {\left(\frac{\mathrm{pe}}{\mathrm{pc}}\right)}^{\frac{\gamma-1}{\gamma}}}\]

\[\fitmath{\mathrm{Te}=(3400)\cdot {\left(\frac{(60000)}{(5000000)}\right)}^{\frac{(1.2)-1}{(1.2)}}\approx 1626.831\,\mathrm{K}}\]

\item Velocidad de salida.

\[\fitmath{\mathrm{Vex}=\sqrt{\frac{2\cdot \gamma}{\gamma-1}\cdot R_g\cdot \mathrm{Tc}\cdot \left(1-{\left(\frac{\mathrm{pe}}{\mathrm{pc}}\right)}^{\frac{\gamma-1}{\gamma}}\right)}}\]

\[\fitmath{\mathrm{Vex}=\sqrt{\frac{2\cdot (1.2)}{(1.2)-1}\cdot (380)\cdot (3400)\cdot \left(1-{\left(\frac{(60000)}{(5000000)}\right)}^{\frac{(1.2)-1}{(1.2)}}\right)}\approx 2843.528\,\mathrm{m/s}}\]

\item Densidad de salida.

\[\fitmath{\mathrm{rhoe}=\frac{\mathrm{pe}}{R_g\cdot \mathrm{Te}}}\]

\[\fitmath{\mathrm{rhoe}=\frac{(60000)}{(380)\cdot (1626.83107)}\approx 0.09705663\,\mathrm{kg/m^{3}}}\]

\item Área de salida por continuidad.

\[\fitmath{\mathrm{Ae}=\frac{\mathrm{mdotc}}{\mathrm{rhoe}\cdot \mathrm{Vex}}}\]

\[\fitmath{\mathrm{Ae}=\frac{(114.1115471)}{(0.09705662733)\cdot (2843.527795)}\approx 0.4134728\,\mathrm{m^{2}}}\]

\item Relación de áreas de la tobera.

\[\fitmath{\mathrm{eps}=\frac{\mathrm{Ae}}{\mathrm{At}}}\]

\[\fitmath{\mathrm{eps}=\frac{(0.4134727761)}{(0.04)}\approx 10.33682\,}\]

\item Empuje al nivel del mar.

\[\fitmath{\mathrm{Fsl}=\mathrm{mdotc}\cdot \mathrm{Vex}+\left(\mathrm{pe}-\mathrm{pa}\right)\cdot \mathrm{Ae}}\]

\[\fitmath{\mathrm{Fsl}=(114.1115471)\cdot (2843.527795)+\left((60000)-(101325)\right)\cdot (0.4134727761)\approx 307392.6\,\mathrm{N}}\]

\item Empuje en vacío.

\[\fitmath{\mathrm{Fvac}=\mathrm{mdotc}\cdot \mathrm{Vex}+\mathrm{pe}\cdot \mathrm{Ae}}\]

\[\fitmath{\mathrm{Fvac}=(114.1115471)\cdot (2843.527795)+(60000)\cdot (0.4134727761)\approx 349287.7\,\mathrm{N}}\]

\item Impulso específico al nivel del mar.

\[\fitmath{\mathrm{Ispsl}=\frac{\mathrm{Fsl}}{\mathrm{mdotc}\cdot 9.80665}}\]

\[\fitmath{\mathrm{Ispsl}=\frac{(307392.5935)}{(114.1115471)\cdot 9.80665}\approx 274.6902\,\mathrm{s}}\]

\item Impulso específico en vacío.

\[\fitmath{\mathrm{Ispvac}=\frac{\mathrm{Fvac}}{\mathrm{mdotc}\cdot 9.80665}}\]

\[\fitmath{\mathrm{Ispvac}=\frac{(349287.7225)}{(114.1115471)\cdot 9.80665}\approx 312.1282\,\mathrm{s}}\]

\item Mach de salida.

\[\fitmath{\mathrm{Mex}=\frac{\mathrm{Vex}}{\sqrt{\gamma\cdot R_g\cdot \mathrm{Te}}}}\]

\[\fitmath{\mathrm{Mex}=\frac{(2843.527795)}{\sqrt{(1.2)\cdot (380)\cdot (1626.83107)}}\approx 3.301443\,}\]

\end{enumerate}


\section{Simulacro 6}
\textbf{Cuestiones.} Un corte longitudinal soporta la presión sobre el área $2rL$ con dos secciones resistentes $2tL$, de donde $\sigma_\theta=pr/t$; un corte transversal soporta $p\pi r^2$ con sección $2\pi rt$, de donde $\sigma_x=pr/(2t)$. La fórmula de Euler es válida si la esbeltez supera la límite $\pi\sqrt{E/\sigma_y}$; al duplicar la longitud la carga crítica se divide por cuatro, y con empotramiento doble se multiplica por cuatro (longitud efectiva $L/2$). La velocidad de maniobra es la mayor a la que un mando brusco completo no supera el factor de carga límite: $V_A=V_s\sqrt{n_{\max}}$.

\textbf{Fuselaje.} A $2400\,\mathrm m$ la temperatura es $272.55\,\mathrm K$ y $p_{cab}=75625.7\,\mathrm{Pa}$. A $11000\,\mathrm m$ la presión es $22632\,\mathrm{Pa}$, de modo que $\Delta p=52993.6\,\mathrm{Pa}$. Resulta $\sigma_\theta=6.29299\times10^{7}\,\mathrm{Pa}$, $\sigma_x=3.1465\times10^{7}\,\mathrm{Pa}$ y $\sigma_{vM}=5.44989\times10^{7}\,\mathrm{Pa}$, con un coeficiente de seguridad estático $\sigma_y/\sigma_{vM}=5.5047$. Ese valor no cubre la fatiga por ciclos de presurización, ni grietas, ni concentraciones en ventanas y uniones; para ello se analiza la vida con curvas $S$-$N$ y la propagación de grietas con criterio de tolerancia al daño.

\textbf{Pandeo.} $I=4.32157\times10^{-8}\,\mathrm{m^4}$, $A=0.000238761\,\mathrm{m^2}$ y $P_{cr}=\pi^2EI/L^2=20733.7\,\mathrm N$. El radio de giro es $0.0134536\,\mathrm m$ y la esbeltez $89.1953$, mayor que la límite $47.9886$, de modo que Euler es aplicable. La carga de fluencia es $71628.3\,\mathrm N$, mayor que $P_{cr}$, así que gobierna el pandeo. Con empotramiento doble la carga crítica sería cuatro veces mayor.

\textbf{Maniobra y ráfaga.} $V_s=\sqrt{2W/(\rho_0SC_{L\max})}=28.3323\,\mathrm{m/s}$ y $V_A=V_s\sqrt{3.8}=55.2299\,\mathrm{m/s}$. El incremento por ráfaga es $\Delta n=\rho_0U_{de}Va/(2W/S)=1.42017$ y el factor resultante $2.42017$, inferior a $3.8$. El modelo ignora el factor de atenuación de ráfaga y la respuesta dinámica del ala, y supone sustentación lineal con la incidencia.

\subsection*{Procedimiento y operaciones del simulacro 6}
La presión de cabina y la exterior salen de ISA; el recipiente delgado se resuelve con equilibrio de fuerzas en cada dirección. La barra se estudia con Euler y se compara con la carga de fluencia. La maniobra sigue $n=(V/V_s)^2$ y la ráfaga discreta se aproxima con $\Delta n=\rho_0U_{de}Va/(2W/S)$.

\begin{enumerate}

\item Temperatura ISA a la altitud de cabina.

\[\fitmath{\mathrm{Tcab}=288.15-0.0065\cdot 2400\approx 272.55\,\mathrm{K}}\]

\item Presión de cabina a 2400 m.

\[\fitmath{\mathrm{pcab}=101325\cdot {\left(\frac{\mathrm{Tcab}}{288.15}\right)}^{\frac{9.80665}{287.05287\cdot 0.0065}}}\]

\[\fitmath{\mathrm{pcab}=101325\cdot {\left(\frac{(272.55)}{288.15}\right)}^{\frac{9.80665}{287.05287\cdot 0.0065}}\approx 75625.66\,\mathrm{Pa}}\]

\item Temperatura ISA en la base de la estratosfera.

\[\fitmath{T_{11}=288.15-0.0065\cdot 11000\approx 216.65\,\mathrm{K}}\]

\item Presión exterior de crucero a 11000 m.

\[\fitmath{\mathrm{pext}=101325\cdot {\left(\frac{T_{11}}{288.15}\right)}^{\frac{9.80665}{287.05287\cdot 0.0065}}}\]

\[\fitmath{\mathrm{pext}=101325\cdot {\left(\frac{(216.65)}{288.15}\right)}^{\frac{9.80665}{287.05287\cdot 0.0065}}\approx 22632.04\,\mathrm{Pa}}\]

\item Diferencia de presión que soporta el fuselaje.

\[\fitmath{\mathrm{dp}=\mathrm{pcab}-\mathrm{pext}}\]

\[\fitmath{\mathrm{dp}=(75625.66139)-(22632.0401)\approx 52993.62\,\mathrm{Pa}}\]

\item Tensión circunferencial de recipiente cilíndrico de pared delgada.

\[\fitmath{\mathrm{sth}=\frac{\mathrm{dp}\cdot 1.9}{0.0016}}\]

\[\fitmath{\mathrm{sth}=\frac{(52993.6213)\cdot 1.9}{0.0016}\approx 6.292993\times10^{7}\,\mathrm{Pa}}\]

\item Tensión longitudinal.

\[\fitmath{\mathrm{sx}=\frac{\mathrm{sth}}{2}}\]

\[\fitmath{\mathrm{sx}=\frac{(62929925.29)}{2}\approx 3.146496\times10^{7}\,\mathrm{Pa}}\]

\item Límite elástico docente.

\[\fitmath{\mathrm{sy}=300000000\approx 3\times10^{8}\,\mathrm{Pa}}\]

\item Tensión equivalente de von Mises biaxial.

\[\fitmath{\mathrm{svm}=\sqrt{{\mathrm{sth}}^{2}-\mathrm{sth}\cdot \mathrm{sx}+{\mathrm{sx}}^{2}}}\]

\[\fitmath{\mathrm{svm}=\sqrt{{(62929925.29)}^{2}-(62929925.29)\cdot (31464962.65)+{(31464962.65)}^{2}}\approx 5.449891\times10^{7}\,\mathrm{Pa}}\]

\item Coeficiente de seguridad estático frente a fluencia.

\[\fitmath{\mathrm{SF}=\frac{\mathrm{sy}}{\mathrm{svm}}}\]

\[\fitmath{\mathrm{SF}=\frac{(300000000)}{(54498913.96)}\approx 5.504697\,}\]

\item Momento de inercia del tubo.

\[\fitmath{I=\frac{\pi}{64}\cdot \left({0.04}^{4}-{0.036}^{4}\right)\approx 4.321575\times10^{-8}\,\mathrm{m^{4}}}\]

\item Área de la sección anular.

\[\fitmath{A=\frac{\pi}{4}\cdot \left({0.04}^{2}-{0.036}^{2}\right)\approx 0.000238761\,\mathrm{m^{2}}}\]

\item Módulo de Young docente.

\[\fitmath{E=7\times10^{10}\approx 7\times10^{10}\,\mathrm{Pa}}\]

\item Carga crítica de Euler biarticulada.

\[\fitmath{\mathrm{Pcr}=\frac{{\pi}^{2}\cdot E\cdot I}{{1.2}^{2}}}\]

\[\fitmath{\mathrm{Pcr}=\frac{{\pi}^{2}\cdot (7\times10^{10})\cdot (4.321574854\times10^{-8})}{{1.2}^{2}}\approx 20733.72\,\mathrm{N}}\]

\item Carga de fluencia a compresión.

\[\fitmath{\mathrm{Py}=\mathrm{sy}\cdot A}\]

\[\fitmath{\mathrm{Py}=(300000000)\cdot (0.0002387610417)\approx 71628.31\,\mathrm{N}}\]

\item Radio de giro de la sección.

\[\fitmath{\mathrm{rg}=\sqrt{\frac{I}{A}}}\]

\[\fitmath{\mathrm{rg}=\sqrt{\frac{(4.321574854\times10^{-8})}{(0.0002387610417)}}\approx 0.01345362\,\mathrm{m}}\]

\item Esbeltez.

\[\fitmath{\mathrm{lam}=\frac{1.2}{\mathrm{rg}}}\]

\[\fitmath{\mathrm{lam}=\frac{1.2}{(0.01345362405)}\approx 89.1953\,}\]

\item Esbeltez límite de validez de la fórmula de Euler.

\[\fitmath{\mathrm{lamc}=\pi\cdot \sqrt{\frac{E}{\mathrm{sy}}}}\]

\[\fitmath{\mathrm{lamc}=\pi\cdot \sqrt{\frac{(7\times10^{10})}{(300000000)}}\approx 47.98862\,}\]

\item Peso de la aeronave.

\[\fitmath{W=11800\approx 11800\,\mathrm{N}}\]

\item Superficie alar.

\[\fitmath{S=16\approx 16\,\mathrm{m^{2}}}\]

\item Velocidad de pérdida a nivel del mar con n=1.

\[\fitmath{\mathrm{Vs}=\sqrt{\frac{2\cdot W}{1.225\cdot S\cdot 1.5}}}\]

\[\fitmath{\mathrm{Vs}=\sqrt{\frac{2\cdot (11800)}{1.225\cdot (16)\cdot 1.5}}\approx 28.33233\,\mathrm{m/s}}\]

\item Velocidad de maniobra para el factor de carga límite.

\[\fitmath{\mathrm{VA}=\mathrm{Vs}\cdot \sqrt{3.8}}\]

\[\fitmath{\mathrm{VA}=(28.33233292)\cdot \sqrt{3.8}\approx 55.22988\,\mathrm{m/s}}\]

\item Incremento del factor de carga por ráfaga discreta sin atenuación.

\[\fitmath{\mathrm{dn}=\frac{1.225\cdot 45\cdot 5\cdot 7.6}{2\cdot \frac{W}{S}}}\]

\[\fitmath{\mathrm{dn}=\frac{1.225\cdot 45\cdot 5\cdot 7.6}{2\cdot \frac{(11800)}{(16)}}\approx 1.420169\,}\]

\item Factor de carga resultante en la ráfaga.

\[\fitmath{\mathrm{nR}=1+\mathrm{dn}}\]

\[\fitmath{\mathrm{nR}=1+(1.420169492)\approx 2.420169\,}\]

\end{enumerate}


\section{Simulacro 7}
\textbf{Cuestiones.} La velocidad indicada es la que mide el anemómetro a partir de la presión dinámica; la verdadera la corrige por densidad, y la velocidad respecto al suelo suma el viento a la verdadera. En una esfera, la ruta más corta es un arco de círculo máximo, que corta a los meridianos con ángulos distintos, de modo que el rumbo varía; la loxodrómica mantiene rumbo constante y se usa en tramos cortos o cerca del ecuador. Con densidad menor aumentan la velocidad de pérdida y de despegue y disminuyen el empuje disponible, por lo que se alarga el recorrido.

\textbf{Vientos.} El viento de $030^\circ$ forma $60^\circ$ con la derrota. La componente transversal es $12.9904\,\mathrm{m/s}$ y la longitudinal, de cara, $7.5\,\mathrm{m/s}$. La corrección de deriva es $\mathrm{WCA}=\arcsin(\mathrm{cross}/\mathrm{TAS})=9.97422^\circ$; el viento llega por la izquierda de la derrota y empuja hacia la derecha, de modo que el rumbo es $090^\circ-9.97422^\circ\approx080.0^\circ$. La velocidad respecto al suelo es $\mathrm{GS}=\mathrm{TAS}\cos\mathrm{WCA}-\mathrm{head}=66.3664\,\mathrm{m/s}$ y el tramo de $180\,\mathrm{km}$ dura $45.2036\,\mathrm{min}$.

\textbf{Ortodrómica.} Con haversine, $0.00143582$ es la haversine del ángulo central, que vale $0.0758026\,\mathrm{rad}$, y la distancia $482.938\,\mathrm{km}$. El rumbo inicial es $\arctan(\mathrm{by}/\mathrm{bx})=77.147^\circ$; como el denominador es positivo no hace falta corregir el cuadrante. El rumbo cambia a lo largo de la ruta porque los meridianos convergen hacia el polo.

\textbf{Despegue.} Al nivel del mar, $V_s=30.9295\,\mathrm{m/s}$, $V_{LOF}=37.1154\,\mathrm{m/s}$, $a=1.79789\,\mathrm{m/s^2}$ y el recorrido es $383.103\,\mathrm m$. A $2000\,\mathrm m$ con ISA$+20$, $T=295.15\,\mathrm K$, $p=79495.2\,\mathrm{Pa}$ y $\rho=0.938288\,\mathrm{kg/m^3}$; entonces $V_s=35.3405\,\mathrm{m/s}$, $V_{LOF}=42.4086\,\mathrm{m/s}$, $T=2451.04\,\mathrm N$, $a=1.30823\,\mathrm{m/s^2}$ y el recorrido es $687.373\,\mathrm m$, es decir $1.79423$ veces el anterior. El modelo omite la resistencia aerodinámica y la sustentación durante la rodadura, el viento, la pendiente de pista y la variación real del empuje con la velocidad.

\subsection*{Procedimiento y operaciones del simulacro 7}
El triángulo de vientos se resuelve proyectando el viento sobre la derrota; la corrección de deriva sale del seno del ángulo transversal. La distancia ortodrómica usa la fórmula haversine y el rumbo inicial la ecuación de rumbo esférico. El despegue se reduce a movimiento uniformemente acelerado hasta $V_{LOF}=1.2V_s$ con la densidad de cada caso.

\begin{enumerate}

\item Velocidad verdadera.

\[\fitmath{\mathrm{TAS}=75\approx 75\,\mathrm{m/s}}\]

\item Intensidad del viento.

\[\fitmath{\mathrm{Vw}=15\approx 15\,\mathrm{m/s}}\]

\item Ángulo entre la procedencia del viento y la derrota.

\[\fitmath{\mathrm{ang}=\frac{\pi\cdot 60}{180}\approx 1.047198\,\mathrm{rad}}\]

\item Componente transversal del viento.

\[\fitmath{\mathrm{cross}=\mathrm{Vw}\cdot \sin\left(\mathrm{ang}\right)}\]

\[\fitmath{\mathrm{cross}=(15)\cdot \sin\left((1.047197551)\right)\approx 12.99038\,\mathrm{m/s}}\]

\item Componente longitudinal del viento, de cara.

\[\fitmath{\mathrm{head}=\mathrm{Vw}\cdot \cos\left(\mathrm{ang}\right)}\]

\[\fitmath{\mathrm{head}=(15)\cdot \cos\left((1.047197551)\right)\approx 7.5\,\mathrm{m/s}}\]

\item Ángulo de corrección de deriva.

\[\fitmath{\mathrm{WCA}=\arcsin\left(\frac{\mathrm{cross}}{\mathrm{TAS}}\right)}\]

\[\fitmath{\mathrm{WCA}=\arcsin\left(\frac{(12.99038106)}{(75)}\right)\approx 0.174083\,\mathrm{rad}}\]

\item Ángulo de deriva en grados.

\[\fitmath{\mathrm{WCAdeg}=\frac{\mathrm{WCA}\cdot 180}{\pi}}\]

\[\fitmath{\mathrm{WCAdeg}=\frac{(0.1740830106)\cdot 180}{\pi}\approx 9.974222\,{}^{\circ}}\]

\item Velocidad respecto al suelo.

\[\fitmath{\mathrm{GS}=\mathrm{TAS}\cdot \cos\left(\mathrm{WCA}\right)-\mathrm{head}}\]

\[\fitmath{\mathrm{GS}=(75)\cdot \cos\left((0.1740830106)\right)-(7.5)\approx 66.36643\,\mathrm{m/s}}\]

\item Tiempo del tramo de 180 km.

\[\fitmath{\mathrm{tmin}=\frac{\frac{180000}{\mathrm{GS}}}{60}}\]

\[\fitmath{\mathrm{tmin}=\frac{\frac{180000}{(66.36643351)}}{60}\approx 45.20357\,\mathrm{min}}\]

\item Latitud del origen.

\[\fitmath{\mathrm{ph}_{1}=\frac{\pi\cdot 40.47}{180}\approx 0.7063347\,\mathrm{rad}}\]

\item Latitud del destino.

\[\fitmath{\mathrm{ph}_{2}=\frac{\pi\cdot 41.3}{180}\approx 0.720821\,\mathrm{rad}}\]

\item Diferencia de longitud, de oeste a este.

\[\fitmath{\mathrm{dlam}=\frac{\pi\cdot \left(2.08+3.56\right)}{180}\approx 0.09843657\,\mathrm{rad}}\]

\item Haversine del ángulo central.

\[\fitmath{\mathrm{hav}={\sin\left(\frac{\mathrm{ph}_{2}-\mathrm{ph}_{1}}{2}\right)}^{2}+\cos\left(\mathrm{ph}_{1}\right)\cdot \cos\left(\mathrm{ph}_{2}\right)\cdot {\sin\left(\frac{\mathrm{dlam}}{2}\right)}^{2}}\]

\[\fitmath{\mathrm{hav}={\sin\left(\frac{(0.7208209811)-(0.7063347483)}{2}\right)}^{2}+\cos\left((0.7063347483)\right)\cdot \cos\left((0.7208209811)\right)\cdot {\sin\left(\frac{(0.09843656981)}{2}\right)}^{2}\approx 0.00143582\,}\]

\item Ángulo central.

\[\fitmath{\mathrm{cang}=2\cdot \arcsin\left(\sqrt{\mathrm{hav}}\right)}\]

\[\fitmath{\mathrm{cang}=2\cdot \arcsin\left(\sqrt{(0.001435819614)}\right)\approx 0.07580257\,\mathrm{rad}}\]

\item Distancia ortodrómica con radio medio.

\[\fitmath{\mathrm{dist}=6371\cdot \mathrm{cang}}\]

\[\fitmath{\mathrm{dist}=6371\cdot (0.07580256811)\approx 482.9382\,\mathrm{km}}\]

\item Numerador de la ecuación del rumbo inicial.

\[\fitmath{\mathrm{by}=\sin\left(\mathrm{dlam}\right)\cdot \cos\left(\mathrm{ph}_{2}\right)}\]

\[\fitmath{\mathrm{by}=\sin\left((0.09843656981)\right)\cdot \cos\left((0.7208209811)\right)\approx 0.07383249\,}\]

\item Denominador de la ecuación del rumbo inicial.

\[\fitmath{\mathrm{bx}=\cos\left(\mathrm{ph}_{1}\right)\cdot \sin\left(\mathrm{ph}_{2}\right)-\sin\left(\mathrm{ph}_{1}\right)\cdot \cos\left(\mathrm{ph}_{2}\right)\cdot \cos\left(\mathrm{dlam}\right)}\]

\[\fitmath{\mathrm{bx}=\cos\left((0.7063347483)\right)\cdot \sin\left((0.7208209811)\right)-\sin\left((0.7063347483)\right)\cdot \cos\left((0.7208209811)\right)\cdot \cos\left((0.09843656981)\right)\approx 0.01684622\,}\]

\item Rumbo inicial, válido porque el denominador es positivo.

\[\fitmath{\mathrm{brg}=\frac{\arctan\left(\frac{\mathrm{by}}{\mathrm{bx}}\right)\cdot 180}{\pi}}\]

\[\fitmath{\mathrm{brg}=\frac{\arctan\left(\frac{(0.07383249289)}{(0.01684622024)}\right)\cdot 180}{\pi}\approx 77.14696\,{}^{\circ}}\]

\item Peso al despegue.

\[\fitmath{\mathrm{Wt}=15000\approx 15000\,\mathrm{N}}\]

\item Superficie alar.

\[\fitmath{\mathrm{Sw}=16\approx 16\,\mathrm{m^{2}}}\]

\item Temperatura ISA a 2000 m.

\[\fitmath{\mathrm{Tisa}=288.15-0.0065\cdot 2000\approx 275.15\,\mathrm{K}}\]

\item Presión ISA a 2000 m.

\[\fitmath{p2k=101325\cdot {\left(\frac{\mathrm{Tisa}}{288.15}\right)}^{\frac{9.80665}{287.05287\cdot 0.0065}}}\]

\[\fitmath{p2k=101325\cdot {\left(\frac{(275.15)}{288.15}\right)}^{\frac{9.80665}{287.05287\cdot 0.0065}}\approx 79495.2\,\mathrm{Pa}}\]

\item Temperatura ISA+20.

\[\fitmath{\mathrm{Thot}=\mathrm{Tisa}+20}\]

\[\fitmath{\mathrm{Thot}=(275.15)+20\approx 295.15\,\mathrm{K}}\]

\item Densidad del día cálido.

\[\fitmath{\mathrm{rhoh}=\frac{p2k}{287.05287\cdot \mathrm{Thot}}}\]

\[\fitmath{\mathrm{rhoh}=\frac{(79495.20193)}{287.05287\cdot (295.15)}\approx 0.9382882\,\mathrm{kg/m^{3}}}\]

\item Velocidad de pérdida a nivel del mar.

\[\fitmath{\mathrm{Vs}_{1}=\sqrt{\frac{2\cdot \mathrm{Wt}}{1.225\cdot \mathrm{Sw}\cdot 1.6}}}\]

\[\fitmath{\mathrm{Vs}_{1}=\sqrt{\frac{2\cdot (15000)}{1.225\cdot (16)\cdot 1.6}}\approx 30.92948\,\mathrm{m/s}}\]

\item Velocidad de despegue a nivel del mar.

\[\fitmath{\mathrm{Vlof}_{1}=1.2\cdot \mathrm{Vs}_{1}}\]

\[\fitmath{\mathrm{Vlof}_{1}=1.2\cdot (30.92947871)\approx 37.11537\,\mathrm{m/s}}\]

\item Aceleración media con empuje constante y rozamiento.

\[\fitmath{a_{1}=9.80665\cdot \left(\frac{3200}{\mathrm{Wt}}-0.03\right)}\]

\[\fitmath{a_{1}=9.80665\cdot \left(\frac{3200}{(15000)}-0.03\right)\approx 1.797886\,\mathrm{m/s^{2}}}\]

\item Recorrido de despegue a nivel del mar.

\[\fitmath{s_{1}=\frac{{\mathrm{Vlof}_{1}}^{2}}{2\cdot a_{1}}}\]

\[\fitmath{s_{1}=\frac{{(37.11537445)}^{2}}{2\cdot (1.797885833)}\approx 383.103\,\mathrm{m}}\]

\item Velocidad de pérdida en el día cálido.

\[\fitmath{\mathrm{Vs}_{2}=\sqrt{\frac{2\cdot \mathrm{Wt}}{\mathrm{rhoh}\cdot \mathrm{Sw}\cdot 1.6}}}\]

\[\fitmath{\mathrm{Vs}_{2}=\sqrt{\frac{2\cdot (15000)}{(0.9382881596)\cdot (16)\cdot 1.6}}\approx 35.34049\,\mathrm{m/s}}\]

\item Velocidad de despegue.

\[\fitmath{\mathrm{Vlof}_{2}=1.2\cdot \mathrm{Vs}_{2}}\]

\[\fitmath{\mathrm{Vlof}_{2}=1.2\cdot (35.34048674)\approx 42.40858\,\mathrm{m/s}}\]

\item Empuje reducido proporcionalmente a la densidad.

\[\fitmath{T_{2}=\frac{3200\cdot \mathrm{rhoh}}{1.225}}\]

\[\fitmath{T_{2}=\frac{3200\cdot (0.9382881596)}{1.225}\approx 2451.038\,\mathrm{N}}\]

\item Aceleración media.

\[\fitmath{a_{2}=9.80665\cdot \left(\frac{T_{2}}{\mathrm{Wt}}-0.03\right)}\]

\[\fitmath{a_{2}=9.80665\cdot \left(\frac{(2451.038458)}{(15000)}-0.03\right)\approx 1.308232\,\mathrm{m/s^{2}}}\]

\item Recorrido de despegue en el día cálido.

\[\fitmath{s_{2}=\frac{{\mathrm{Vlof}_{2}}^{2}}{2\cdot a_{2}}}\]

\[\fitmath{s_{2}=\frac{{(42.40858409)}^{2}}{2\cdot (1.308232253)}\approx 687.3734\,\mathrm{m}}\]

\item Cociente de recorridos.

\[\fitmath{\mathrm{ratio}=\frac{s_{2}}{s_{1}}}\]

\[\fitmath{\mathrm{ratio}=\frac{(687.3733623)}{(383.1030299)}\approx 1.794226\,}\]

\end{enumerate}


\section{Simulacro 8}
\textbf{Cuestiones.} En retracción el aceite actúa sobre el área anular, pistón menos vástago: la misma presión da menos fuerza y el mismo caudal da más velocidad. La capacidad nominal es la corriente por hora a un régimen de descarga y temperatura concretos; la utilizable es menor por profundidad de descarga, envejecimiento y temperatura. La ecuación de Breguet integra $dW/W$ suponiendo $V$, $L/D$ y consumo constantes; si cambian con la masa o la altitud hay que integrar por tramos o usar una variante con ascenso escalonado.

\textbf{Actuador.} Las áreas son $A_p=0.0019635\,\mathrm{m^2}$ y $A_r=0.000490874\,\mathrm{m^2}$. Las fuerzas son $F_{ext}=41233.4\,\mathrm N$ y $F_{ret}=30925.1\,\mathrm N$. Con el caudal dado, $v_{ext}=0.509296\,\mathrm{m/s}$ y $v_{ret}=0.679061\,\mathrm{m/s}$. Con $30\,\mathrm{kN}$ de carga, que es inferior a $F_{ext}$, la potencia hidráulica es $21000\,\mathrm W$, la útil $15278.9\,\mathrm W$ y la disipada $5721.13\,\mathrm W$ en válvulas y fugas. La bomba toma $24705.9\,\mathrm W$. Sin enfriamiento, el fluido se calienta a $0.167284\,\mathrm{K/s}$, es decir, unos $10\,\mathrm K$ por minuto, lo que justifica un intercambiador de calor.

\textbf{Batería.} La autonomía es $0.8\cdot40/15=2.13333\,\mathrm h$. El cable tiene $R=\rho L/A=0.0344\,\Omega$, con caída $0.516\,\mathrm V$ y pérdida $7.74\,\mathrm W$. La tensión en la carga es $23.484\,\mathrm V$ si se toman $24\,\mathrm V$ en bornes de la batería, una caída del $2.15\,\%$.

\textbf{Breguet.} Con $c_w=c\,g_0=0.000166713\,\mathrm{s^{-1}}$, el alcance es $R=(V/c_w)(L/D)\ln(W_0/W_1)=4.92564\times10^{6}\,\mathrm m\approx4926\,\mathrm{km}$ y la autonomía $E=5.94884\,\mathrm h$. La fracción de combustible es $1-W_1/W_0=0.2$. La ecuación supone empuje igual a resistencia, vuelo horizontal y velocidad y consumo constantes, y no incluye reservas, ascenso ni descenso.

\subsection*{Procedimiento y operaciones del simulacro 8}
En el actuador se aplican $F=pA$ y $v=Q/A$ con el área útil de cada cámara; las potencias salen del balance hidráulico. La batería se trata como reserva energética utilizable y el cable como una resistencia en serie. El alcance sigue la integración de Breguet para reactor con $L/D$ y consumo constantes.

\begin{enumerate}

\item Presión del sistema.

\[\fitmath{p=21000000\approx 2.1\times10^{7}\,\mathrm{Pa}}\]

\item Área del pistón.

\[\fitmath{\mathrm{Ap}=\frac{\pi}{4}\cdot {0.05}^{2}\approx 0.001963495\,\mathrm{m^{2}}}\]

\item Área del vástago.

\[\fitmath{\mathrm{Ar}=\frac{\pi}{4}\cdot {0.025}^{2}\approx 0.0004908739\,\mathrm{m^{2}}}\]

\item Fuerza máxima en extensión.

\[\fitmath{\mathrm{Fext}=p\cdot \mathrm{Ap}}\]

\[\fitmath{\mathrm{Fext}=(21000000)\cdot (0.001963495408)\approx 41233.4\,\mathrm{N}}\]

\item Fuerza máxima en retracción, con área anular.

\[\fitmath{\mathrm{Fret}=p\cdot \left(\mathrm{Ap}-\mathrm{Ar}\right)}\]

\[\fitmath{\mathrm{Fret}=(21000000)\cdot \left((0.001963495408)-(0.0004908738521)\right)\approx 30925.05\,\mathrm{N}}\]

\item Caudal.

\[\fitmath{Q=0.001\approx 0.001\,\mathrm{m^{3}/s}}\]

\item Velocidad de extensión.

\[\fitmath{\mathrm{vext}=\frac{Q}{\mathrm{Ap}}}\]

\[\fitmath{\mathrm{vext}=\frac{(0.001)}{(0.001963495408)}\approx 0.5092958\,\mathrm{m/s}}\]

\item Velocidad de retracción para el mismo caudal.

\[\fitmath{\mathrm{vret}=\frac{Q}{\mathrm{Ap}-\mathrm{Ar}}}\]

\[\fitmath{\mathrm{vret}=\frac{(0.001)}{(0.001963495408)-(0.0004908738521)}\approx 0.6790611\,\mathrm{m/s}}\]

\item Potencia hidráulica entregada al actuador.

\[\fitmath{\mathrm{Ph}=p\cdot Q}\]

\[\fitmath{\mathrm{Ph}=(21000000)\cdot (0.001)\approx 21000\,\mathrm{W}}\]

\item Potencia mecánica útil con carga de 30 kN.

\[\fitmath{\mathrm{Pu}=30000\cdot \mathrm{vext}}\]

\[\fitmath{\mathrm{Pu}=30000\cdot (0.5092958179)\approx 15278.87\,\mathrm{W}}\]

\item Potencia disipada como calor en válvulas y fugas.

\[\fitmath{\mathrm{loss}=\mathrm{Ph}-\mathrm{Pu}}\]

\[\fitmath{\mathrm{loss}=(21000)-(15278.87454)\approx 5721.125\,\mathrm{W}}\]

\item Potencia de entrada a la bomba con rendimiento 0.85.

\[\fitmath{\mathrm{Pin}=\frac{\mathrm{Ph}}{0.85}}\]

\[\fitmath{\mathrm{Pin}=\frac{(21000)}{0.85}\approx 24705.88\,\mathrm{W}}\]

\item Calentamiento si 18 kg de fluido absorben el calor sin enfriar.

\[\fitmath{\mathrm{dTdt}=\frac{\mathrm{loss}}{18\cdot 1900}}\]

\[\fitmath{\mathrm{dTdt}=\frac{(5721.125463)}{18\cdot 1900}\approx 0.1672844\,\mathrm{K/s}}\]

\item Autonomía con 80 % de la capacidad utilizable.

\[\fitmath{t=\frac{0.8\cdot 40}{15}\approx 2.133333\,\mathrm{h}}\]

\item Resistencia del cableado de cobre de ida y vuelta.

\[\fitmath{R=\frac{1.72\times10^{-8}\cdot 12}{6\times10^{-6}}\approx 0.0344\,\mathrm{ohm}}\]

\item Caída de tensión en el cable.

\[\fitmath{\mathrm{dV}=15\cdot R}\]

\[\fitmath{\mathrm{dV}=15\cdot (0.0344)\approx 0.516\,\mathrm{V}}\]

\item Pérdida por efecto Joule en el cable.

\[\fitmath{\mathrm{Ploss}={15}^{2}\cdot R}\]

\[\fitmath{\mathrm{Ploss}={15}^{2}\cdot (0.0344)\approx 7.74\,\mathrm{W}}\]

\item Tensión en la carga.

\[\fitmath{\mathrm{Vload}=24-\mathrm{dV}}\]

\[\fitmath{\mathrm{Vload}=24-(0.516)\approx 23.484\,\mathrm{V}}\]

\item Caída relativa en por ciento.

\[\fitmath{\mathrm{frac}=\frac{\mathrm{dV}}{24}\cdot 100}\]

\[\fitmath{\mathrm{frac}=\frac{(0.516)}{24}\cdot 100\approx 2.15\,}\]

\item Velocidad de crucero.

\[\fitmath{V=230\approx 230\,\mathrm{m/s}}\]

\item Consumo específico por unidad de peso de empuje.

\[\fitmath{\mathrm{cw}=1.7\times10^{-5}\cdot 9.80665\approx 0.000166713\,\mathrm{1/s}}\]

\item Eficiencia aerodinámica.

\[\fitmath{\mathrm{LD}=16\approx 16\,}\]

\item Logaritmo de la razón de pesos inicial y final.

\[\fitmath{\mathrm{lnW}=\ln\left(1.25\right)\approx 0.2231436\,}\]

\item Alcance de Breguet.

\[\fitmath{\mathrm{Rbr}=\frac{V}{\mathrm{cw}}\cdot \mathrm{LD}\cdot \mathrm{lnW}}\]

\[\fitmath{\mathrm{Rbr}=\frac{(230)}{(0.00016671305)}\cdot (16)\cdot (0.2231435513)\approx 4925639\,\mathrm{m}}\]

\item Autonomía de Breguet.

\[\fitmath{\mathrm{Eh}=\frac{\frac{\mathrm{LD}\cdot \mathrm{lnW}}{\mathrm{cw}}}{3600}}\]

\[\fitmath{\mathrm{Eh}=\frac{\frac{(16)\cdot (0.2231435513)}{(0.00016671305)}}{3600}\approx 5.948839\,\mathrm{h}}\]

\item Fracción de combustible respecto al peso inicial.

\[\fitmath{\mathrm{ff}=1-\frac{1}{1.25}\approx 0.2\,}\]

\end{enumerate}

